\section{Cumplimiento de los módulos UX/UI exigidos}
\label{sec:modulos}

El enunciado exige los \textbf{cuatro} módulos obligatorios y \textbf{tres
de los cuatro} módulos opcionales. Esta sección integra la respuesta del
proveedor Claude (prompt 04), seleccionada por su mayor precisión
mecanística (mapeo concreto de variables del aura, descripción del
\textit{pipeline} de roles del LLM, justificación estructurada del
módulo descartado) y por su trazabilidad explícita con el backlog de
RF/RNF de la sección~\ref{sec:requerimientos}.

\subsection{Módulos obligatorios}

\subsubsection{1. Interfaz gráfica no convencional}
El simulador implementa una sala virtual tridimensional renderizada con
WebGL/Three.js, en la que el entrevistador es un avatar estilizado
rodeado por un \textbf{aura luminosa procedural} que actúa como canal de
feedback continuo. El aura mapea las métricas multimodales a parámetros
visuales independientes:

\begin{itemize}
  \item \textbf{Color}: codifica la valencia global del desempeño
        (ámbar a verde).
  \item \textbf{Intensidad}: codifica el nivel de fluidez verbal.
  \item \textbf{Ritmo de pulsación}: codifica la cadencia detectada.
  \item \textbf{Asimetría espacial}: codifica el equilibrio entre canal
        verbal y no verbal.
\end{itemize}

\noindent La interacción ocurre dentro del entorno tridimensional sin
ventanas modales sobreimpuestas, lo que constituye una desviación del
paradigma de formulario web convencional. La visualización mantiene
$\ge 30$\,fps en hardware de referencia y degrada a una representación
bidimensional simplificada cuando se detecta GPU insuficiente.
\textbf{Cubre} \rf{05}, \rf{15}, \rnf{03}.

\subsubsection{2. Interfaz de voz}
Canal de voz bidireccional cuasi-real:

\begin{itemize}
  \item Las respuestas del entrevistador LLM se sintetizan mediante TTS
        (modelo de Gemini con voz seleccionable).
  \item Las intervenciones del candidato se transcriben mediante STT
        continuo con detección de fin de habla y segmentación por
        turnos.
  \item Sobre la transcripción y la señal acústica se calculan
        \textbf{en cliente} las métricas prosódicas que alimentan el
        aura (fluidez, ritmo, pausas, tono, densidad de muletillas).
  \item Una \textbf{capa de comandos de voz} permite controlar la
        sesión mediante intenciones reservadas (``pausar'',
        ``repetir pregunta'', ``siguiente'', ``terminar entrevista'')
        procesadas por un reconocedor de gramática restringida
        independiente del canal sustantivo.
\end{itemize}

\noindent La latencia mediana de TTS y de aparición de subtítulos se
mantiene por debajo del umbral fijado por \rnf{01}--\rnf{02}.
\textbf{Cubre} \rf{02}, \rf{06}, \rf{17}--\rf{19}, \rnf{01}--\rnf{02}.

\subsubsection{3. Integración con uno o varios LLM}
El sistema integra \textbf{tres roles diferenciados} sobre la API de
Gemini, cada uno con su propio prompt sistémico, contexto y memoria
acotada:

\begin{itemize}
  \item \textbf{Entrevistador}: genera preguntas adaptadas a industria,
        rol y nivel del candidato; reformula en función del historial
        conversacional de la sesión.
  \item \textbf{Coach}: opera en modo asincrónico al cierre de la
        sesión; consume la transcripción anonimizada y las métricas
        agregadas para producir un plan de mejora estructurado por
        competencias.
  \item \textbf{Moderador}: actúa exclusivamente en sesiones de peer
        mock; observa la transcripción cruzada de ambos pares y emite
        intervenciones breves cuando detecta desviaciones del
        protocolo o silencios prolongados.
\end{itemize}

\noindent Ningún rol recibe el video crudo del candidato: recibe
únicamente los descriptores numéricos derivados del \textit{pipeline}
de MediaPipe en cliente. \textbf{Cubre} \rf{02}, \rf{07}, \rf{16},
\rnf{05}, \rnf{14}.

\subsubsection{4. Ayuda contextual para el usuario}
Cuatro capas observables:

\begin{enumerate}
  \item \textbf{Tour guiado} al primer ingreso, navegable con teclado,
        que destaca cada zona funcional de la sala virtual.
  \item \textbf{Etiquetas emergentes} (\textit{tooltips}) accesibles
        por foco y por puntero sobre todos los controles, con
        descripción breve y atajo de teclado asociado.
  \item \textbf{Asistente coach-on-demand}, invocable mediante el
        comando ``ayuda'' o tecla dedicada, que responde dudas sobre
        el funcionamiento del simulador (no sobre el contenido
        evaluativo de la entrevista) consultando una base documental
        indexada.
  \item \textbf{Inferencia de necesidad de ayuda} por inactividad o
        errores repetidos: si el usuario permanece sin acción durante
        un umbral configurable o reintenta el mismo comando con falla,
        el sistema sugiere la sección de ayuda relevante.
\end{enumerate}

\noindent \textbf{Cubre} \rf{15}--\rf{16}.

\subsection{Módulos opcionales seleccionados (3 de 4)}

\subsubsection{Personalización continua}
\begin{itemize}
  \item Historial conversacional persistente, consultable como línea
        de tiempo de sesiones con métricas, transcripción y reporte
        (\rf{08}).
  \item Plan adaptativo recalculado al cierre de cada sesión: el
        módulo \textit{Coach} prioriza las competencias con menor
        desempeño en las últimas $N$ sesiones y selecciona los
        siguientes ejercicios en función de esa priorización (\rf{07}).
  \item Perfil de industria y rol declarado por el usuario y refinado
        a partir del CV cargado y de las respuestas previas; condiciona
        el banco de preguntas y el vocabulario técnico esperado por el
        entrevistador (\rf{01}).
  \item \textbf{Calibración progresiva del nivel} mediante un modelo
        de estimación de habilidad por competencia (similar a IRT) que
        sube o baja la dificultad de la siguiente pregunta según la
        calidad de la respuesta previa, con nivel inicial fijado por
        una sesión diagnóstica (\rf{09}).
  \item Persistencia de preferencias de interfaz, accesibilidad y
        configuración del aura entre sesiones y dispositivos, con
        sincronización por cuenta autenticada (\rnf{04}).
\end{itemize}

\subsubsection{Interactividad con otros usuarios}
\begin{itemize}
  \item Peer mock vía WebRTC con señalización por servidor propio y
        media P2P entre pares; sala con dos roles activos
        (candidato y observador) y un moderador LLM compartido
        (\rf{10}).
  \item Panel del observador con métricas multimodales del candidato
        en vivo (con consentimiento explícito previo del candidato
        registrado en sesión) y formulario estructurado para emitir
        feedback por competencia (\rf{11}).
  \item Comentarios anclados al \textit{timestamp} de la grabación
        local de la sesión, navegables por marcador y editables por el
        observador hasta el cierre (\rf{11}).
  \item Intercambio de roles mediante un único control que invierte
        candidato y observador y reinicia el contador de la nueva
        entrevista, preservando el historial de la sesión anterior
        como adjunto consultable (\rf{12}).
  \item \textbf{Modo panel} con hasta tres observadores simultáneos
        sobre un único candidato, con feedback agregado al cierre y
        resolución de conflictos por promedio ponderado por
        antigüedad del observador.
  \item \textbf{Feedback cruzado asincrónico}: solicitar revisión de
        una entrevista grabada por un par seleccionado del directorio,
        con plazo configurable y notificación al solicitante cuando el
        feedback queda disponible.
\end{itemize}

\subsubsection{Gamificación}
\begin{itemize}
  \item Rangos por competencia (comunicación verbal, lenguaje
        corporal, contenido técnico, manejo del estrés) calculados
        como agregaciones móviles sobre las últimas $N$ sesiones, con
        umbrales fijos y observables en el perfil (\rf{13}).
  \item \textit{Badges} otorgados por logros verificables y
        deterministas (p.\,ej. completar diez sesiones consecutivas,
        alcanzar puntaje umbral en una competencia, completar un peer
        mock como observador), con criterios documentados y reversibles
        ante anulación de la sesión correspondiente (\rf{13}).
  \item Ligas temáticas agrupadas por industria y rol, con
        clasificación semanal basada en \emph{mejora relativa} del
        usuario respecto a su propia línea base, no en comparación
        absoluta entre usuarios, para evitar penalizar a principiantes
        (\rf{14}).
  \item Retos grupales con objetivo común (acumulación colectiva de
        horas de práctica o de peer mocks completados en un periodo) y
        progreso visible para todos los miembros del grupo.
\end{itemize}

\subsection{Módulo opcional descartado: interfaces hápticas}

Se descarta la inclusión de interfaces hápticas en el alcance del MVP
por cuatro razones convergentes:

\begin{enumerate}
  \item \textbf{Dominio}: la entrevista laboral simulada no requiere
        intercambio táctil; la totalidad del valor evaluativo y
        formativo se transmite por canales auditivo, visual y textual,
        y no existe evidencia documentada de mejora de desempeño por
        háptica en este dominio.
  \item \textbf{Hardware}: la única superficie háptica realmente
        accesible en el contexto de uso esperado es la vibración del
        móvil, un canal de muy baja resolución expresiva que además no
        está disponible en escritorio (entorno principal del simulador).
  \item \textbf{Inmersión}: una vibración durante la entrevista
        interrumpiría la simulación de un contexto formal y se
        percibiría como notificación intrusiva en lugar de información
        integrada, contradiciendo la metáfora del
        \textit{Warachikuy}\footnote{El \textit{Warachikuy} (del quechua
        \textit{wara}, taparrabo otorgado en la mayoría de edad) era el
        rito de iniciación a la adultez en el Tahuantinsuyu, en el que
        los jóvenes superaban pruebas rigurosas de valor y destreza
        para recibir el reconocimiento de la comunidad y de la
        autoridad. La ceremonia, declarada de interés cultural por el
        Ministerio de Cultura del Perú, se reescenifica anualmente en
        la explanada de Saqsayhuamán. El proyecto adopta esta metáfora
        andina —en lugar de la del \textit{dojo} oriental utilizada en
        la propuesta inicial— porque captura con mayor precisión el
        carácter de tránsito profesional, examen estructurado y
        reconocimiento al cierre que define a una entrevista laboral.}.
  \item \textbf{Recursos}: el equipo prefiere profundizar los tres
        módulos opcionales seleccionados, donde existen requerimientos
        automatizables y verificables, antes que dispersar esfuerzo en
        una capa háptica cuya verificación automática sería difícil.
\end{enumerate}

\noindent La decisión queda sujeta a revisión en una fase posterior si
se incorpora un caso de uso específico (por ejemplo, alertas hápticas
para usuarios con discapacidad auditiva durante peer mock) que
justifique el costo de implementación.
